\documentclass[10pt,letterpaper]{article}
\usepackage{technical-implementation}
\setDocTitle{Laboratory 05: Trust Boundary Threat Model}
\setDocSubtitle{Controlled security engineering and verification}
\setDocVersion{1.0}
\setDocStatus{Independent engineering guidance}
\setDocOwner{latex-docs / security}
\setDocDate{2026-11-02}
\setDocSummary{Independent reproducible laboratory with negative-path acceptance criteria.}
\begin{document}
\makedoctocpage
\begin{notebox}[Source boundary]
Textbook concepts are anchored to Matt Bishop, \emph{Computer Security: Art and Science}, Second Edition (2019). All procedures, implementation examples, scoring rules and modern framework mappings below are independently authored supplemental engineering material, not excerpts, official compliance guidance or evidence of deployed controls.
\end{notebox}
\section{Objective and book anchors}
Read Bishop Chapters Trust Boundary Threat Model before execution. This laboratory independently examines trust boundary threat model. It is not a reproduction of the book's exercises.
\section{Prerequisites}
1, 4, 14. All tests run on synthetic data in an isolated, disposable environment. Never use real credentials, customer records, malware or third-party targets.
\section{Reproducible procedure}
\begin{enumerate}
\item Identify scope, source commit and evaluator. Write the property, threat actor, trust assumption and expected reject behavior.
\item Review \texttt{Paper or PlantUML; isolated document API model}. Make a baseline hypothesis before executing anything.
\item Execute or adapt the documented command and record stdout, stderr, exit code, compiler/interpreter version and timestamp.
\item Run a permitted case, a prohibited case and a failure/boundary case. Use the corresponding security evidence form.
\item Repeat after a controlled defect or policy modification; observe whether the negative test would catch the regression.
\item Return the environment to its original state and submit the evidence for a second-person review.
\end{enumerate}
\section{Command}
\begin{viewbox}\small\ttfamily procedures/threat-modeling-procedure.tex\end{viewbox}
\section{Required engineering exercise}
Record assets and entry points in the Threat Model Template; create an abuse-case per trust boundary.
\section{Safety and assurance limits}
Separate attacker action from vulnerable assumption and business impact; record evidence for a proposed control.
\section{Acceptance criteria}
\begin{enumerate}
\item Hypothesis and preconditions are explicit and the test procedure is reproducible.
\item Correct authorization or safety behavior is observable on positive and negative paths.
\item Source revision, relevant tool versions, observed result and open uncertainties are recorded.
\item A reviewer can independently reproduce the result without production access.
\end{enumerate}
\section{Evidence package}
Include artifact ID, control objective, negative test ID, environment manifest, results, residual risk, reviewer and disposition. Evaluation applies only to this modeled and executed scope.

\end{document}
